\documentclass[11pt]{article}
\usepackage[T1]{fontenc}
\usepackage[margin=1in]{geometry}
\usepackage{amsmath,amssymb,amsthm,mathtools}
\usepackage{graphicx}
\usepackage{microtype}
\usepackage[hidelinks]{hyperref}

\newtheorem{theorem}{Theorem}
\newtheorem{corollary}[theorem]{Corollary}
\newcommand{\A}{\mathcal A}
\newcommand{\Sset}{\mathcal S}
\newcommand{\Tc}{\mathcal T_{\!c}}
\newcommand{\T}{\mathcal T}

\title{An exact fixed-point formula for the\\consistent Bellman operator}
\author{Automated Math Discovery Run \texttt{fa\_banach\_001}\\Model: GPT5.6}
\date{22 July 2026}

\begin{document}
\maketitle

\begin{center}
\fbox{\parbox{0.89\linewidth}{\textbf{Claimed status: full quantitative
answer, subject to human review.}  The formula below completely determines
the consistent Bellman fixed point from the classical optimal $Q$-function
and the one-step self-transition probabilities.}}
\end{center}

\section{Source question}

Kadurha David asks how the unique fixed point of the consistent Bellman
operator relates to the fixed point obtained from the classical Bellman
operator \cite[Chapter 5, Section 5.1.2, printed p.~34]{source}.  The passage
appears on PDF page 38:

\begin{figure}[h]
  \centering
  \includegraphics[width=0.96\linewidth]{figures/open_problem_crop.png}
  \caption{The open question in the source paper.}
\end{figure}

Bellemare et al. had already proved the qualitative facts that the consistent
operator is optimality-preserving and gap-increasing
\cite[Corollary 3]{bellemare}.  The result here gives the missing exact
coordinatewise relation.

\section{Setting and notation}

Let $\Sset$ and $\A$ be nonempty finite sets, let
$P(\,cdot\mid s,a)$ be transition probabilities, let
$r:\Sset\times\A\to\mathbb R$, and let $0\leq\gamma<1$.  For a
$Q$-function $q$, put $v_q(s)=\max_{a\in\A}q(s,a)$.  The classical and
consistent Bellman operators are
\begin{align*}
 (\T q)(s,a)
   &=r(s,a)+\gamma\sum_{x\in\Sset}P(x\mid s,a)v_q(x),\\
 (\Tc q)(s,a)
   &=r(s,a)+\gamma\left(
       P(s\mid s,a)q(s,a)
       +\sum_{x\neq s}P(x\mid s,a)v_q(x)
     \right).
\end{align*}
Both are $\gamma$-contractions in the supremum norm.  Write $q^*$ for the
fixed point of $\T$, $q_c$ for the fixed point of $\Tc$, and
\[
  v^*(s)=\max_a q^*(s,a),
  \qquad d(s,a)=P(s\mid s,a).
\]

\section{Proof intuition}

At a consistent fixed point, the term caused by returning immediately to the
same state contains $q_c(s,a)$ itself.  Moving that term to the left eliminates
the self-loop and produces an operator on state values alone.  Its contraction
factor is no larger than $\gamma$, even after division by
$1-\gamma d(s,a)$.  The classical optimal value $v^*$ is a fixed point of
this new operator: every action gives at most $v^*(s)$, and a classical
optimal action gives equality.  Uniqueness therefore forces
$\max_a q_c(s,a)=v^*(s)$.  Substitution then yields the exact formula, not
just an inequality between action gaps.

\section{Exact relation}

\begin{theorem}[Coordinatewise fixed-point formula]
For every $(s,a)\in\Sset\times\A$,
\begin{equation}\label{eq:main}
 \boxed{
 v^*(s)-q_c(s,a)
   =\frac{v^*(s)-q^*(s,a)}{1-\gamma d(s,a)}.}
\end{equation}
Equivalently,
\begin{equation}\label{eq:qformula}
 q_c(s,a)
 =\frac{q^*(s,a)-\gamma d(s,a)v^*(s)}
        {1-\gamma d(s,a)}.
\end{equation}
In particular, $\max_a q_c(s,a)=v^*(s)$ for every state.
\end{theorem}

\begin{proof}
For a bounded function $v:\Sset\to\mathbb R$, define
\begin{align*}
 H_a v(s)
   &=\frac{r(s,a)+\gamma\sum_{x\neq s}P(x\mid s,a)v(x)}
           {1-\gamma d(s,a)},\\
 Hv(s)&=\max_{a\in\A}H_a v(s).
\end{align*}
The denominator is positive.  For each $s,a$,
\begin{align*}
 |H_a v(s)-H_a w(s)|
 &\leq
 \frac{\gamma(1-d(s,a))}{1-\gamma d(s,a)}\,
 \|v-w\|_\infty\\
 &\leq \gamma\|v-w\|_\infty,
\end{align*}
because $1-d\leq 1-\gamma d$.  Taking maxima does not increase this bound,
so $H$ is a $\gamma$-contraction and has a unique fixed point.

Put $v_c(s)=\max_a q_c(s,a)$.  The equation $\Tc q_c=q_c$, with its
self-loop term moved to the left, says exactly
\[
 q_c(s,a)=H_a v_c(s).
\]
Taking the maximum over $a$ gives $v_c=Hv_c$.

We next show that $v^*$ is also an $H$-fixed point.  The classical fixed-point
equation gives
\[
 q^*(s,a)=r(s,a)+\gamma d(s,a)v^*(s)
      +\gamma\sum_{x\neq s}P(x\mid s,a)v^*(x).
\]
Consequently,
\begin{equation}\label{eq:keyidentity}
 v^*(s)-H_a v^*(s)
 =\frac{v^*(s)-q^*(s,a)}{1-\gamma d(s,a)}\geq0.
\end{equation}
Since $\A$ is finite, some action $a^*$ attains
$q^*(s,a^*)=v^*(s)$; for this action equality holds in
\eqref{eq:keyidentity}.  Hence $Hv^*=v^*$.  Uniqueness of the $H$-fixed point
now yields $v_c=v^*$.

Finally, $q_c(s,a)=H_a v^*(s)$.  Replacing the numerator of $H_a v^*$ by
$q^*(s,a)-\gamma d(s,a)v^*(s)$ proves \eqref{eq:qformula}; subtracting it
from $v^*(s)$ gives \eqref{eq:main}.
\end{proof}

\begin{corollary}[Exact policy and gap consequences]
For every state $s$:
\begin{enumerate}
\item $\operatorname*{arg\,max}_a q_c(s,a)
      =\operatorname*{arg\,max}_a q^*(s,a)$, including all ties;
\item each action gap is amplified by the factor
      $(1-\gamma d(s,a))^{-1}\in[1,(1-\gamma)^{-1}]$;
\item $q_c(s,a)\leq q^*(s,a)$, and the inequality is strict exactly when
      $q^*(s,a)<v^*(s)$ and $d(s,a)>0$.
\end{enumerate}
\end{corollary}

\begin{proof}
The first two assertions follow immediately from \eqref{eq:main}.  For the
third, subtract \eqref{eq:qformula} from $q^*(s,a)$ to obtain
\[
 q^*(s,a)-q_c(s,a)
 =\frac{\gamma d(s,a)}{1-\gamma d(s,a)}
   \bigl(v^*(s)-q^*(s,a)\bigr).
\]
\end{proof}

\section{Verification and review notes}

The script \texttt{code/check\_random\_mdps.py} generated 400 random finite
MDPs (two to seven states, two to five actions), mixed in self-loop weights up
to $0.995$, and independently iterated both operators.  The worst observed
direct formula error was $1.08\times10^{-11}$; the worst residual after
substituting the formula into $\Tc$ was $2.07\times10^{-13}$; and the maximum
state values agreed to displayed floating-point precision.  These checks are
not part of the proof.

The theorem is deliberately stated for finite tabular MDPs, matching the
source's summation formulas and attained maxima.  The proof extends to broader
bounded settings under the usual measurability and supremum hypotheses, but
that extension is not claimed here.  The main human-review point is the
reduction from the $Q$-fixed-point equation to the contraction $H$.

\section{Bounded novelty check}

The run's registry, solution, attempt, and proof-gap indexes had no hit for
arXiv:2410.03706 or this fixed-point relation.  Searches over arXiv, AAAI,
NeurIPS, and OpenReview combined ``consistent Bellman operator'' with
``fixed point'', ``formula'', ``self transition'', and ``action gap''.
The primary qualitative predecessor \cite{bellemare} and later reviews were
inspected.  They establish optimality preservation and gap increase, but the
explicit coordinate identity \eqref{eq:main} was not found.  This was a
bounded, not exhaustive, search; novelty confidence is therefore moderate and
is limited to the exact formula and its sharp consequences.

\begin{thebibliography}{9}
\bibitem{source}
K. Kadurha David,
\emph{Topological Foundations of Reinforcement Learning},
arXiv:2410.03706, 2024.
\url{https://arxiv.org/abs/2410.03706}

\bibitem{bellemare}
M. G. Bellemare, G. Ostrovski, A. Guez, P. S. Thomas, and R. Munos,
\emph{Increasing the Action Gap: New Operators for Reinforcement Learning},
Proc. AAAI Conf. Artificial Intelligence 30 (2016), 1476--1483;
arXiv:1512.04860.
\url{https://doi.org/10.1609/aaai.v30i1.10303}
\end{thebibliography}

\end{document}

